\begin{textsample}{POS Dim 1 – vem\_ed\_22 – Score 3.00 – vem\_ed\_22/t111.txt}  \label{ex:f1_pos_176}
PCIs em 2023: resultados e indicadores O ano de 2023 marcou a projeção internacional dos PCIs / Cesu, como se verá em detalhe ao \textbf{longo} das próximas páginas.
Expandir as colaborações e manter a qualidade dos projetos \textbf{foram} metas perseguidas e alcançadas: ocorreram 109 PCIs no ano, dos quais 27 novos (25 \%) e 82 reeditados (75 \%), com a participação de 91 professores de Fatecs, 78 de instituições internacionais, 2.505 estudantes estrangeiros e 2.450 alunos de Fatecs.
De 2013 a 2023, cerca de 10 mil fatecanos (9.934) participaram de mais de 500 PCIs.
Tais experiências \textbf{foram} amplamente \textbf{compartilhadas} em comunicações acadêmicas: em 2023, a equipe dos PCIs participou de 12 eventos nacionais, 11 internacionais, 10 grupos especializados, acompanhou 8 docentes estrangeiros em visitas a Fatecs, \textbf{publicou} 2 artigos, apresentou 32 trabalhos em congressos nacionais e 31 em eventos internacionais.
Além disso, organizou o Simpósio Internacional IVES Cesu 2023 e a Jornada de PCIs, noticiados, respectivamente, em VEm 18 e VEm 20.
Os principais resultados e indicadores estão sintetizados no infográfico a \textbf{seguir}.

% matched lemmas: compartilhar, longo, publicar, seguir, ser
\end{textsample}
